\documentclass[output=paper,modfonts,nonflat,colorlinks,citecolor=brown]{langscibook}
\ChapterDOI{10.5281/zenodo.18234081}
\author{Bruno Olsson\affiliation{University of Regensburg}}
\title{Negation in Kashaya}

\abstract{Negation in Kashaya, a Native North American language, is described based on a textual corpus. The main strategies for clausal negation are a negative verb suffix and a negative clitic, and I show that their distribution can be predicted to a large extent from the inflectional form of the verb. Negation of nominal and adjectival predicates, as well as constituent negation, employs the negative clitic. A special negative existential verb is employed to negate existential and locative predications. I pay special attention to a non-negative use of the negative clitic, which is used to form (positive) indefinite pronouns, a phenomenon that, to my knowledge, has not been described for any other languages.
\keywords{Pomoan languages, negation, evidentiality, switch-reference, indefinite pronouns}
}

\IfFileExists{../localcommands.tex}{
   % add all extra packages you need to load to this file

\usepackage{tabularx,multicol}
\usepackage{url}
\urlstyle{same}

\usepackage{listings}
\lstset{basicstyle=\ttfamily,tabsize=2,breaklines=true}

% \usepackage{langsci-basic}
\usepackage{langsci-optional}
\usepackage{langsci-lgr}
\usepackage{langsci-gb4e}

\usepackage[linguistics,edges]{forest}
\usepackage{subfigure}

\usepackage{arydshln}

\usepackage{pifont}
\newcommand{\cmark}{\ding{51}}

\usepackage{tabto}


\usepackage{multirow} %for Tsova-Tush, NGT, panara
\usepackage{booktabs} %for Tsova-Tush, panara
\usepackage{graphicx} %for Tsova-Tush
\usepackage{longtable} %for Tsova-Tush, NGT
\usepackage{vcell} %for Tsova-Tush

% \usepackage{pdflscape} %kambaata -- has been added to enable landscape formatting of table 3
\usepackage{xcolor} %kambaata,ngarinjin,nivkh -- has been added to enable gray shading in tables
\definecolor{light-gray}{gray}{0.8}%ngarinjin
\definecolor{lightest-gray}{gray}{0.95}%ngarinjin

%\usepackage{subcaption} %for NGT, but does not work with subfigure package from above
\usepackage{slgloss} %for NGT
% \newfontfamily{\handshape}{handshape.TTF} %for NGT %HC2022-11-04: it does not work and creates an error "Package fontspec Error: The font "handshape" cannot be found", eventhough I have uploaded the file handshape.TTF, so I outcommented it for the time being. Could you solve this Sebastian? thanks :)
%\newcommand\customfont[1]{{\usefont{T1}{handshape}{m}{n} #1 }} %for NGT
\newfontfamily\ToneFont{CharisSIL-Regular.ttf}
\AdditionalFontImprint{Charis SIL}
\usepackage[shortlabels]{enumitem} %NGT 


\usepackage{leipzig} %for the glosses in panara
\usepackage{glossaries} %panara
% \usepackage{lscape} %panara
% \input{libertine-preamble} %panara
% \setmathfont{Libertinus Math} %panara
%   \usepackage{enumerate} %panara - it clashes with enumitem which is really needed for NGT, and it works in panara also with enumitem, i.e. we don't need enumerate.

% \setmainfont[Ligatures=TeX,Numbers=OldStyle]{Linux Libertine O} %panara
% \setsansfont[Ligatures=TeX,Numbers=OldStyle]{Linux Biolinum O} %panara

\usepackage{hyperref}%kalamang,yukana
% \usepackage{tipx}%kalamang,ungarinjin, yukana

% \usepackage{wrapfig}%ungarinjin
% \usepackage{tabulary}%ungarinjin
%\usepackage{subfig}%ungarinjin %deleted bc subfigure already above
\PassOptionsToPackage{colorinlistoftodos,disable}{todonotes}
\usepackage{todonotes}
\usepackage{hyphenat}
\usepackage{langsci-branding}


   %\newcommand*{\orcid}{}

\makeatletter
\let\thetitle\@title
\let\theauthor\@author
\makeatother

\newcommand{\togglepaper}[1][0]{
   \bibliography{../localbibliography}
   \papernote{\scriptsize\normalfont
     \theauthor.
     \titleTemp.
     To appear in:
     Matti Miestamo \& Ljuba Veselinova (ed.).
     Title of book.
     Berlin: Language Science Press. [preliminary page numbering]
   }
   \pagenumbering{roman}
   \setcounter{chapter}{#1}
   \addtocounter{chapter}{-1}
}


\renewbibmacro*{index:name}[5]{% %Tsova-Tush
  \usebibmacro{index:entry}{#1} %Tsova-Tush
    {\iffieldundef{usera}{}{\thefield{usera}\actualoperator}\mkbibindexname{#2}{#3}{#4}{#5}}} %for Tsova-Tush

    
\newcommand{\gl}[1]{\textsc{#1}}

\newcommand{\ipa}[1]{{\textit{#1}}} %geshiza 
\newcommand{\ipaEX}[1]{{\textit{#1}}} %geshiza


%\newcommand{\hspaceThis}[1]{\hphantom{#1}} %this was in localcommands in the kambaata chapter

\newcommand{\ANA}{\textsc{ana}} %gyeli
\newcommand{\ATTR}{\textsc{attr}} %gyeli
\newcommand{\CONJ}{\textsc{conj}} %gyeli
\newcommand{\ID}{\textsc{id}} %gyeli
\newcommand{\IDEO}{\textsc{ideo}} %gyeli
\newcommand{\INCH}{\textsc{inch}} %gyeli
\newcommand{\LINK}{\textsc{link}} %gyeli
\newcommand{\NP}{NP} %gyeli
\newcommand{\NCA}{\textsc{nca}} %gyeli
\newcommand{\NPI}{\textsc{npi}} %gyeli
\newcommand{\PCF}{\textsc{pcf}} %gyeli
\newcommand{\PN}{\textsc{pn}} %gyeli
\newcommand{\PREP}{\textsc{prep}} %gyeli
\newcommand{\PRIOR}{\textsc{prior}} %gyeli
\newcommand{\PRO}{\textsc{pro}} %gyeli
\newcommand{\PROSP}{\textsc{prosp}} %gyeli
\newcommand{\R}{\textsc{r}} %gyeli
\newcommand{\RETRO}{\textsc{retro}} %gyeli
\newcommand{\SEQU}{\textsc{sequ}} %gyeli
\newcommand{\STAMP}{\textsc{stamp}} %gyeli
\newcommand{\SUB}{\textsc{sub}} %gyeli
\newcommand{\TM}{TM} %gyeli
\newcommand{\V}{\textsc{v}} %gyeli


\newcommand{\hspaceThis}[1]{\hphantom{#1}} %for Kambaata chapter, helps for \db (which allows alignment of word and gloss and not without taking into account the "s)


%for Panara:
\renewbibmacro*{index:name}[5]{%
  \usebibmacro{index:entry}{#1}
    {\iffieldundef{usera}{}{\thefield{usera}\actualoperator}\mkbibindexname{#2}{#3}{#4}{#5}}}

 %panara
\newleipzig{ade}{ade}{adessive} %panara
\newleipzig{dir}{dir}{directional} %panara
\newleipzig{emph}{emph}{emphasis} %panara
\newleipzig{fnl}{fnl}{final} %panara
\newleipzig{nml}{nml}{nominal form}%panara
\newleipzig{ine}{ine}{inessive}  %panara 
\newleipzig{ints}{ints}{intensifier} %panara
\newleipzig{iter}{iter}{iterative} %panara
\newleipzig{plac}{plac}{pluractional} %panara


% \makeglossaries %panara %HC not sure what to do with these following lines, but it does not work yet. I have added the abbreviations in the same table is in other chapters, so this glossary function is not needed anymore.
% \newglossarystyle{myglosses}{%
%   \renewenvironment{theglossary}%
% {\begin{multicols}{2}\raggedright}
% {\end{multicols}}
%     \renewcommand*{\glossaryheader}{}
%     \renewcommand*{\glsgroupheading}[1]{}
%     \renewcommand*{\glsgroupskip}{}
%     \renewcommand*{\glsclearpage}{} 
% % set how each entry should appear:
%     \renewcommand*{\glossentry}[2]{
%     \noindent\makebox[7em][l]{\glstarget{##1}{\textsc{\glossentryname{##1}}}}
%         \glossentrydesc{##1}
%     }
%     \renewcommand*{\subglossentry}[3]{%
%         \glossentry{##2}{##3}
%     }
% }

\newcommand{\glopa}{\textsc{dem}} %kalamang
\newcommand{\glme}{\textsc{top}} %kalamang
\newcommand{\glte}{\textsc{nfin}} %kalamang
\newcommand{\glse}{\textsc{iam}} %kalamang
\newcommand{\glkin}{\textsc{vol}} %kalamang
\newcommand{\glteba}{\textsc{prog}} %kalamang
\newcommand{\gli}{\textsc{plnk}} %kalamang
\newcommand{\glet}{\textsc{irr}} %kalamang
\newcommand{\glta}{\textit{ta}} %kalamang
\newcommand{\glnn}{\textit{n}} %kalamang

\newcommand{\tabitem}{~~\llap{\textbullet}~~}%kogi

%For Matti and Ljuba's comments on the last version of the chapters:

\newcommand{\matti}[2][color=green!40,size=\small]{\todo[#1]{MM: #2}}
\newcommand{\ljuba}[2][color=blue!40,size=\small]{\todo[#1]{LV: #2}}
\newcommand{\mm}[2][color=green!40,size=\small]{\todo[#1]{MM: #2}}
\renewcommand{\lv}[2][color=blue!40,size=\small]{\todo[#1]{LV: #2}}
\newcommand{\hc}[2][size=\small]{\todo[#1]{HC: #2}}
\newcommand{\gloss}[2][color=red!30,size=\tiny]{\todo[#1]{Glossing: #2}}
\newcommand{\glosst}[2][color=red!30,size=\tiny,inline]{\todo[#1]{Glossing: #2}} %t for (in the) text / inline
\newcommand{\aut}[2][color=yellow!50,size=\tiny]{\todo[#1]{Author: #2}}%for author

\renewcommand{\keywords}[1]{%
  \par\noindent\textbf{Keywords: #1.}%
}


\renewcommand{\REF}[2][]{(\ref{#2}#1)}
\renewcommand{\sectref}[1]{§\ref{#1}}
\newcommand{\Abf}{{Ā}\kern-2pt\raisebox{3.1pt}{́}\kern2pt}
\newcommand{\abf}{{ā}\kern-1pt\raisebox{1.1pt}{́}\kern1pt}
\newcommand{\A}{{Ā}\kern-1.5pt\raisebox{3.1pt}{́}\kern1.5pt}
\newcommand{\U}{{ʊ}\kern-1.5pt\raisebox{-.4pt}{̀}\kern1.5pt}

\newcommand{\fdot}{{f\kern-.75pt\raisebox{-2.5pt}{̣}\kern.75pt}}

\newcommand{\handshape}[2][1.5]{\includegraphics[height=#1ex]{fonts/#2.png}}

    \bibliography{../localbibliography}
    \togglepaper[23]
}{}

\begin{document}
\maketitle

\section{The language}
   
Kashaya (ISO: \href{https://iso639-3.sil.org/code/kju}{kju}, Glottocode: \href{https://glottolog.org/resource/languoid/id/kash1280}{kash1280}) belongs to the Pomoan family, a small family of Northern California consisting of seven closely related languages. The traditional Kashaya territory is in Sonoma County, just north of the San Fransisco Bay. The language is critically endangered with a small number of speakers remaining \parencite[111]{Golla2011}. In this chapter, I describe negation in the language based on the fully segmented and interlinearised corpus made available in \citet{Buckley2019}, a resource based on the 82 texts collected in \citet{Oswalt1964} (and 8 texts from \citealt{Halpern1940}).
The structure of the chapter follows the negation questionnaire designed by the editors of this book (\citetv{chapters/questionnaire}).

The primary research on Kashaya was undertaken in the late 1950s by Robert L. Oswalt, whose grammar (\citealt{Oswalt1961}), text collection (\citealt{Oswalt1964}) and dictionary manuscript (\citealt{Oswalt}) are the main sources on the language. 
Oswalt's grammar focuses on morphophonological processes and morphotactics, and with the exceptions of e.g.\ the evidentials \citep{Oswalt1986} and the switch-reference system \citep{Oswalt1983}, many areas of Kashaya morphosyntax remain unexplored. 
Later research has primarily concerned theoretical aspects of the intricate morphophonology of the language (e.g.\ \citealt{Buckley1994a, Buckley1994}, \citeyear{Buckley1997}).

Kashaya has relatively free word order, although SOV dominates. Verb morphology is mostly suffixing. There is no participant indexing (i.e.\ person/number agreement) on the verb, but there is rich derivational and inflectional encoding of manner, direction, aspect, tense, evidentiality and a range of other categories.

The suffixes are organised according to a system of position classes in \citet{Oswalt1961} and \citet{Buckley1994}. 
One of the position classes, labelled \textsc{Class XIV} in \citet{Oswalt1961}, deserves mention here. This position class contains the largest number of suffixes, and one suffix from this class has to be present in every verb. 
Some common suffixes from this class have evidential and modal meanings. In contexts in which none of these suffixes are appropriate, speakers resort to a default-like suffix that \citet{Oswalt1961} labels the ``Absolutive''. 
The exact meaning of this highly frequent suffix, at least in its uses on main clause verbs, is unclear to me. It is also used to form deverbal nouns, and provides the citation form of verbs. As pointed out by \citet{Walker2016}, Oswalt's term is somewhat confusing for a modern reader, so I opt for Walker's more informative label \textsc{Default Verbal Suffix} (abbreviated \textsc{dvs}).
In the nominal domain, case marking distinguishes semantic agents from patients (cf.\ \sectref{sec:BO:flag}). 
    
    
\section{Clausal negation}

\subsection{Standard negation}\label{sec:BO:standard}
\subsubsection{Overview of the standard negation constructions}
Two options exist for negating declarative verbal main clauses. The first option is to add the negative clitic \textit{=tʰin} to the verb. A declarative main clause is given in (\ref{ex:BO:0}), 
and a negated clause, with the same verb (\textit{ʔdu·cicʼ-} `know') is given in (\ref{ex:BO:1}), illustrating the negative clitic option.%
\footnote{%
An identifier at the start of every example indicates the source: KT for the \textit{Kashaya Texts} \citep{Oswalt1964} and H for the materials in \citet{Halpern1940}; the following numbers indicate text, paragraph and sentence, so `KT 3:19.11' for the 11th sentence in the 19th paragraph in the 3rd of Oswalt's texts. Examples labelled Kashaya Dictionary are taken from \citet{Oswalt}. A preliminary online version of the dictionary is being prepared by Gene Buckley, and can be accessed at \url{https://www.webonary.org/kashaya/}. The first line of the example sentences provides the original orthography. The second line provides a morphemic analysis (based on \citealt{Buckley2019}). Note that clitics (such as the negative \textit{=tʰin}) are written as separate orthographic words in the first line, but indicated by means of equals signs in the second line.
In the Kashaya orthography, <t> is the dental plosive, <ṭ> is the alveolar plosive, <š> is the palatal fricative, <c> is the alveopalatal affricate, and <y> is the palatal approximant. The raised dot, as in <a·>, marks a long vowel. Other symbols have the same values as their IPA equivalents.
} 

{%
\ea\label{ex:BO:0}
\glll {[\ldots]} he· qʼóʔo hlaw ʔdu·ciʔ\\
    {} he· qʼóʔo hlaw ʔdu·cicʼ-ʔ\\
    {} and song also know-\gl{dvs}\\
\glt `\ldots and [they] know songs too.' (KT 47:9.3)\\
\z
}%

{%
\ea\label{ex:BO:1}
\glll {[\ldots]} ʔa mu ʔcayʔ yacoʔkʰe qʼoʔo ṭʼo dú·ciʔ tʰin\\
    {} ʔa mu =ʔcayʔ =yacoʔkʰe qʼoʔo =ʔṭʼo ʔdu·cicʼ-ʔ =tʰin\\
    {} \gl{1sg.nom} that =person =\gl{poss} song =\gl{emph} know-\gl{dvs} =\gl{neg}\\
\glt `[However] I don’t know his song.' ({KT 17:11.2})\\
\z
}%

\noindent The second option is to add the negative suffix \textit{-tʰ} to the verb, as in (\ref{ex:BO:2}) and (\ref{ex:BO:3}).

{%
\ea\label{ex:BO:2}
\glll \llap{``}be·li ṭʼo ma·dú·tʰe·'' nihciʔ\\
    be·li =ʔṭʼo ma·duc-tʰ-a-e· nihcicʼ-ʔ\\
    here =\gl{emph} one.arrive-\gl{neg-vis.ipfv-nfnl} say.\gl{pl.agt-dvs}\\
\glt `“He hasn’t come here, ” they said.' ({KT 58:6.5})\\
\z
}%

\ea\label{ex:BO:3}
\glll \llap{``}ʔacaʔ tʼaʔtʰé· to'' nihcedu\\
    ʔaca·c tʼa·d-tʰ-a-e· to nihced-u\\
    person feel-\gl{neg}-\gl{vis.ipfv-nfnl} \gl{1sg.acc} say-\gl{dvs}\\
\glt `“I didn’t think it was a person" [he said].' ({KT 45:16.4})\\
\z

\noindent The factors that determine the choice between the clitic \textit{=tʰin} and the suffix \textit{-tʰ} are described in \sectref{sec:BO:distr}. The impact of negation on the aspectual value of the verb is described in \sectref{sec:BO:aspect}. The expression of ‘not yet' is described in \sectref{sec:kas:2.1.4}.

\subsubsection{The distribution of \textit{=tʰin} and \textit{-tʰ} as clausal negators}
\label{sec:BO:distr}
The two options differ rather dramatically in frequency in corpus data. \tabref{tab:BO:freqs} shows counts performed on the corpus.
In the 933 instances of negation that were extracted, the suffixal negative \textit{-tʰ} accounted for only 11.1\% of the attestations, while the negative clitic \textit{=tʰin} accounted for the remaining 88.9\%.%
\footnote{This number excludes prohibitives (\sectref{sec:BO:proh}). It also excludes the positive indefinite pronouns, which are formed by adding the negative clitic to question words, as discussed in \sectref{sec:BO:indef}.
}

\begin{table}[t]
\caption{Frequencies of negative structures}
\label{tab:BO:freqs}
 \begin{tabular}{llrrlrr}
  \lsptoprule
        & \multicolumn{3}{l}{Syntactic context} & \multicolumn{3}{l}{Morphological context}\\
  \midrule
Suffix \textit{-tʰ} & Verb & 104 & (100\%) & V\gl{-neg-}EVID & 31 & (29.8\%)\\
 &&&& V\gl{-neg-cond} & 24 &(23.0\%)\\
  &&&& V\gl{-neg-resp} & 16& (15.4\%)\\
 &&&& V\gl{-neg-}SR & 33& (31.7\%)\\ \cmidrule{5-7}
  &&&& & 104& (100\%)\\ \cmidrule{2-7}
  & \bfseries Total: & 104 & (11.1\%)\\
\midrule
Clitic \textit{=tʰin} &  Verb & 644 &(77.6\%) &V\gl{-dvs =neg} & 415 &(64.4\%)\\
&&&& V\gl{-fut =neg} & 219 &(34.0\%)\\
&&&& V\gl{-SR =neg} & 9& (1.3\%)\\
&&&& V\gl{-grl.intv =neg} & 1& (0.2\%)\\ \cmidrule{5-7}
  &&&& & 644& (100\%)\\ \cmidrule{2-7}
    & Non-verb & 185 &(19.8\%)\\
  & \bfseries Total: & 829 & (88.9\%)\\
\midrule
& \bfseries Total: & 933 & (100\%)\\
  \lspbottomrule
 \end{tabular}
 \\
{SR = any of the switch-reference suffixes; EVID = any of the evidential suffixes
}
\end{table}


Both syntactic and morphological context seem to determine the speakers' choice between the suffixal and clitic options. The clitic \textit{=tʰin} accepts any part of speech as its host, whereas the suffix \textit{-tʰ} is restricted to verbal morphology. I extracted 185 instances of \textit{=tʰin} attached to non-verbal hosts (mostly noun phrases; see further \sectref{sec:BO:nonverb} and \sectref{sec:BO:scope}) accounting for about a fifth of its uses. The remaining 644 instances of \textit{=tʰin} attached to verbs, where it is in competition with the suffix \textit{-tʰ}.

The choice between \textit{=tʰin} and \textit{-tʰ} in verbal contexts depends largely on the form of the verb, and more specifically on which of the suffixes in Class XIV is used on the verb. Here I will show that certain suffixes, such as the evidentials and the Conditional \textit{-iʔba}, only occur with suffixal negation, whereas others, such as the Future suffix \textit{-ʔkʰe} and the Default Verbal Suffix \textit{-w}, only occur with the negative clitic. The suffixes marking switch-reference on subordinate (adverbial) clauses occur with both options.

First, I will consider the rich evidential system of Kashaya (\citealt{Oswalt1986}). \tabref{tab:BO:evid} lists suffixes with evidential functions, plus the `Inferential' \textit{-bi}, which seems to have mirative rather than evidential semantics. The labels are self-explanatory except for the `Performatives', which ``signify that the speaker knows of what he speaks because he is performing the act himself or has just performed it'' (\citeyear[34]{Oswalt1986}). %
\footnote{%
The evidential pair \textit{-(w)a} and \textit{-(y)a} are called ``Factual'' and ``Visual'', respectively, in \citet{Oswalt1961} and \citet{Buckley1994}, and ``Visual-Factual'' and ``Visual'' in \citet{Oswalt1990}. I prefer to use the Visual label for both suffixes and distinguish them by their aspectual value.
} %
Note that some evidentials alternate according to the aspectual value of the verb (either imperfective or perfective), which will be relevant for the discussion of aspect and negation in \sectref{sec:BO:aspect}.

\begin{table}[t]
\caption{The Kashaya evidential suffixes}
\label{tab:BO:evid}
 \begin{tabular}{lll}
  \lsptoprule
  & \textsc{ipfv} & \textsc{pfv}\\
  \midrule
Performative  &  \itshape -(w)ela & \itshape -mela  \\
Visual & \itshape -w(a) & \itshape -y(a) \\
Auditory  & \multicolumn{2}{c}{\itshape -inn(a)} \\
Circumstantial  & \multicolumn{2}{c}{\itshape -q(a)} \\
Hearsay\strut & \multicolumn{2}{c}{\itshape -do} \\ \hdashline
Inferential\strut & \multicolumn{2}{c}{\itshape -bi} \\
  \lspbottomrule
 \end{tabular}
\end{table}

The corpus data show that verbs containing an evidential suffix are negated by means of the suffix \textit{-tʰ}, and not by the clitic \textit{=tʰin}. Of the 31 negated verbs with evidential specification in \tabref{tab:BO:freqs}, all exhibit the suffixal option.\footnote{A reviewer correctly points out that the evidential suffixes are not consistently segmented in \citet{Buckley2019}. This problem was easily solved by manually inspecting the 104 verbs with suffixal negation that were extracted.} The relatively low number of negated verbs inflected with evidentials in the corpus makes it difficult to interpret the significance of the association between evidentials and suffixal negation, however, and the grammar (\citealt{Oswalt1961}) does not clarify whether the negative clitic is disallowed or merely dispreferred with verbs marked by evidential suffixes.\footnote{I was able to identify about two dozen additional examples of negated verbs marked for evidentiality in the dictionary manuscript (\citealt{Oswalt}), all of which occurred with the suffixal option.}

Consider next the use of the Default Verbal Suffix. The sparsity of verbs directly inflected for evidentiality finds its correlate in the frequent use of this suffix (which seems to contribute no meaning of its own) in the text corpus. As pointed out by \citet[239]{Oswalt1961}, verbs with the Default Verbal Suffix can only be negated by means of the clitic \textit{=tʰin}. Such verbs consitute the overwhelmingly most common type of negated verb in \tabref{tab:BO:freqs}, with 415 corpus attestations.

There are good reasons to suspect that the prevalence of (negated and non-negated) verbs marked with the Default Verbal Suffix, and the relatively low number of verbs bearing evidential suffixes, is an artefact of the type of materials represented in the corpus. The texts contained in the corpus mostly feature monologic narrative discourse, and it seems that the use of evidential suffixes in this text type is very different from their use in spontaneous, interactive speech. Narrative discourse typically makes use of what \citet[34]{Oswalt1986} refers to as the \textsc{narrative construction}. In this mode, the Hearsay evidential \textit{-do} is suffixed to the Assertive (a defective, auxiliary-like verb), which is added after a demonstrative element (e.g.\ \textit{mul} `that'), and such chunks, which can be glossed as `it is said that', are typically sprinkled throughout the paragraphs, whereas the main verb in each sentence is suffixed with the Default Verbal Suffix. This has the consequence that verbs directly suffixed with evidentials are comparatively rare in the textual corpus (except in reported speech), while verbs with the Default Verbal Suffix are comparatively common. It is therefore likely that a corpus consisting of a larger portion of interactive speech would have contained more verbs directly suffixed with evidentials, and hence more instances of the suffixal negator \textit{-tʰ}.

In the corpus, negated verbs with evidential suffixes are mostly attested in direct speech, as in (\ref{ex:BO:2}) and (\ref{ex:BO:3}) above. It is interesting to note that verbs that are directly suffixed by means of the Hearsay evidential are more common in the older sections of the corpus, in the texts collected by Abraham Halpern in the 1940s. In these texts, one finds examples such as (\ref{ex:BO:Vdo}), with the suffixal negator attached to a verb directly inflected for Hearsay evidentiality.

{%
\ea 
\label{ex:BO:Vdo}
\glll cuʔtʼémtʰido soh ʔacʰuláʔ soh hísʼu· ʔel\\
    cuʔtʼem-tʰ-do soh ʔacʰulaq-ʔ soh hisʼu· =ʔel\\
    shoot-\gl{neg}-\gl{hrsy} just miss.\gl{pl}-\gl{dvs} just arrow =\gl{acc}\\
\glt `But they didn't hit him; the arrows just missed.' ({H 4:3.6})\\
\z
} %

The third domain of verb morphology to consider is future tense marking. The most common markers of future tense are the Future suffix \textit{-ʔkʰe}, which expresses e.g.\ planned, predicted and intended events, as in (\ref{ex:BO:fut1}), and the Performative Intentive suffix \textit{-te}, which only expresses intentions of the speaker, seen on the first verb in (\ref{ex:BO:intent}).

{%
\ea 
\label{ex:BO:fut1}
\glll ʔama hcʰócʼkʰe ʔdo· yal\\
    ʔama· hcʰoc-ʔkʰe ʔ-do-e· yal\\
    land disappear-\gl{fut} \gl{asrt-hrsy-nfnl} \gl{1pl.acc}\\
\glt `We have heard our world is coming to an end.' ({KT 66:2.6})\\
\z
}%

\ea 
\label{ex:BO:intent}
\glll sáhqaʔte· ʔa ʔacaʔ duhkʰuyá·dal ʔa ʔacaʔ bimucí·dal\\
      sahqacʼ-te-e· ʔa ʔaca·c duhkʰuy-ad-a-el ʔa ʔaca·c bimucid-a-el\\
      quit-\gl{prfm.intv}-\gl{nfnl} \gl{1sg.nom} person kill-\gl{dur}-\gl{vis.ipfv}-\gl{acc} \gl{1sg.nom} person eat-\gl{vis.ipfv-acc}\\
\glt `I am going to quit killing people and eating them.' ({KT 6:9.8})\\
\z

\noindent Negated clauses with future time reference seem to be formed only with the Future suffix \textit{-ʔkʰe}, as in (\ref{ex:BO:futneg1}) and (\ref{ex:BO:futneg2}), and never with the Performative Intentive suffix.%
\footnote{
Another suffix, the General Intentive \textit{-ti}, occurs in some subordinate clause types. It occurs negated in the formation of negative purpose clauses, as described in \sectref{sec:BO:NegPurp}.
} %
A verb suffixed with \textit{-ʔkʰe} can only be negated with the clitic \textit{=tʰin}, and not with the suffix \textit{-tʰ} \citep[258]{Oswalt1961}.

\ea 
\label{ex:BO:futneg1}
\glll ʔama hcʰócʼkʰe tʰin [\ldots]\\
    ʔama· hcʰoc-ʔkʰe =tʰin {}\\
    land disappear-\gl{fut} =\gl{neg} {}\\
\glt `The world isn’t going to end [\ldots].' ({KT 66:13.2})\\
\z

{%
\ea 
\label{ex:BO:futneg2}
\glll bihkʰúyʔkʰe tʰin qahmulʔkʰé· ya mito\\
    bihkʰuy-ʔkʰe =tʰin qahmu·l-ʔkʰe-e· ya mito\\
    eat.up-\gl{fut} =\gl{neg} leave.food-\gl{fut-nfnl} \gl{1pl.nom} \gl{2sg.acc}\\
\glt `We won’t eat it up; we’ll save some for you.' ({KT 10:11.5})\\
\z
}%

The Future suffix \textit{-ʔkʰe} also has more broadly modal uses, most prominently in a construction in which it combines with negation to express inability, `cannot', as described in \sectref{sec:BO:inabil}. The Future suffix is very frequent in the corpus, and it is frequent in negated contexts. Its incompatibility with the negative suffix \textit{-tʰ} contributes to the high frequency of the clitic \textit{=tʰin} in \tabref{tab:BO:freqs}.

Two additional suffixes occur only with suffixal negation: the Conditional \mbox{\textit{-iʔba}} and the Responsive \textit{-em}. The Conditional, illustrated in (\ref{ex:BO:cond1}), has various modal meanings relating to ability and epistemic possibility (see also \sectref{sec:BO:SubClauses} and \sectref{sec:BO:inabil}). The uses of the Responsive seem to relate to information structure (cf.\ \citealt[283]{Oswalt1961}; \citealt[32--34]{Olsson2010}), but I leave the functions of this suffix, and its interaction with negation, for future research.

{%
\ea 
\label{ex:BO:cond1}
\glll buṭaqa tʰín he· buṭaqá {pʰiʔyahtʰíʔba, [\ldots]}\\
     buṭaqa =tʰin =he· buṭaqa pʰiʔya·q-tʰ-iʔba\\
     bear =\gl{neg} =or bear recognise.by.sight-\gl{neg}-\gl{cond} \\
\glt `Whether it was a bear or not a bear couldn't be told just by looking, [\ldots].' ({KT 44:2.6})\\
\z
}%

\tabref{tab:BO:Vneg} summarises the choice between suffixal and clitic negators, according to the form of the verb (excluding the Responsive and the General Intentive; for the latter, see \sectref{sec:BO:NegPurp}).
Overall, verbs marked with the Future tense \textit{-ʔkʰe} or the Default Verbal Suffix account for almost all instances of the clitic \textit{=tʰin} in verbal contexts in the corpus. I am not aware of any synchronic explanation for the patterning of the negative clitic with certain verb suffixes, and the negative suffix with others, which seems to be an arbitrary quirk of the grammar. Negation in non-verbal clauses exclusively employs the negative clitic, as will be seen in \sectref{sec:BO:nonverb}.%

Note finally that both the suffixal and clitic options are attested in switch-reference clauses, as discussed further in \sectref{sec:BO:SR}.

 \begin{table}[t]
 \caption{Summary of main types of negated verbs: the suffix \textit{-tʰ} vs.\ the clitic \textit{=tʰin} }
 \label{tab:BO:Vneg}
  \begin{tabular}{lp{2em}p{2em}}
   \lsptoprule
   	& \multicolumn{2}{c}{Negator}\\ \cmidrule{2-3}
  Class XIV-suffix on the verb           & \textit{-tʰ} &  \textit{=tʰin}  \\
   \midrule
 Evidential suffixes & \cmark \\
 Default Verbal Suffix & & \cmark \\
 Future \textit{-ʔkʰe} & & \cmark \\
 Conditional \textit{-iʔba} & \cmark \\
% Responsive \textit{-em} & \cmark \\
 Switch-reference suffixes & \cmark & \cmark \\
   \lspbottomrule
  \end{tabular}
 \end{table}

\subsubsection{Negation and aspect}
\label{sec:BO:aspect}

Aspect is an important component of Kashaya verb morphology. \citet{Oswalt1990} describes the central opposition as one between perfective and imperfective. Many verb roots are inherently perfective or imperfective, others are aspectually neutral. Various derivational suffixes (typically with directional semantics: `up', `away' etc.\@) have perfectivising or imperfectivising functions and can be used to derive a new verb with a different aspectual profile than that of the root (somewhat suggestive of Slavic aspectual systems). Whether a verb form is perfective or imperfective is most clearly seen from some of the evidential suffixes in \tabref{tab:BO:evid}. The Performative and the Visual evidentials occur in pairs in which one member combines with a perfective verb and the other with an imperfective verb (aspectually neutral roots without derivational markers can combine with either member, depending on the meaning that the speaker wants to convey).

But the aspectual potential of verbs is limited when the verb is negated: a verb suffixed with the negative \textit{-tʰ} can only combine with the imperfective members of the Visual and Performative evidentials (\citealt[244--246]{Oswalt1961}). The minimal triplet in (\ref{ex:BO:ipfv}) is taken from \citet[45]{Oswalt1990}. The aspectually neutral root `pack' can in the affirmative combine with either the perfective or imperfective Performative suffixes (\ref{ex:BO:ipfv.a}--\ref{ex:BO:ipfv.b}). Under negation, only the imperfective variant is possible (\ref{ex:BO:ipfv.c}).%
\footnote{
Oswalt does not say explicitly whether (\ref{ex:BO:ipfv.c}) can also have past time reference, i.e.\ `I didn't pack', but this seems to be implied by his discussion of these examples. There are only 6 corpus attestations of the negated Performative, and none of them has unambiguous past time reference, so these attestations shed no light on the issue. Note that (\ref{ex:BO:hide}) below does not have unambiguous past time reference, despite its English translation (`We haven’t hid anybody') which could be rendered in the present without change in meaning (`We aren't hiding anybody.').
} %
This triplet is given without context in the source (and is probably elicited), but I have not found any counterexamples to this pattern in the corpus. The neutralisation of the perfective/imperfective contrast is a clear example of paradigmatic asymmetry under negation (cf.\ \citealt[54]{Miestamo2005}; see also \citealt{MiestamoAuwera2011} for the interaction between negation and aspect cross-lingustically).

\ea
\label{ex:BO:ipfv}
\begin{xlist}
\ex\label{ex:BO:ipfv.a}
\glll qowáhmela\\
     qowaq-mela\\
     pack-\gl{prfm.pfv}\\
\glt `I just packed (a suitcase).'
\ex\label{ex:BO:ipfv.b}
\glll qowá·qala\\
     qowaq-ela\\
     pack-\gl{prfm.ipfv}\\
\glt `I am packing (a suitcase).'
\ex\label{ex:BO:ipfv.c}
\glll qowáhtʰela\\
     qowaq-tʰ-ela\\
     pack-\gl{neg-prfm.ipfv}\\
\glt `I am not packing.'
\end{xlist}
\z

Note that the aspectual membership of the roots (and derivational suffixes) does not affect their combinability with negation. Any type of verb root or derived verb can be negated, but the resulting form is always treated as imperfective by the evidential suffixes that are sensitive to this distinction. It is not clear whether negation by means of the clitic \textit{=tʰin} has any consequences for aspect. The evidential suffixes that reveal the aspectual profile of a verb form do not combine with the negative clitic \textit{=tʰin} (as discussed above), so there is no way of finding this out.

Other than the restriction against the perfective members of those evidential pairs that show a perfective/imperfective contrast, I am not aware of any morphological asymmetries between negative and affirmative verbs.


\subsubsection{Negation and phasal markers}\label{sec:kas:2.1.4}
Many languages employ non-compositional expressions for `not yet'. In Kashaya, the negative continuative `not yet' is expressed compositionally, by simply negating the continuative adverbial \textit{ʔoʔ} `still' (root /ʔoc/), i.e.\ `still not'.
Compare the affirmative clause in (\ref{ex:BO:still}) with the negated clauses in (\ref{ex:BO:notyet1}) and (\ref{ex:BO:notyet2}).

{%
\ea 
\label{ex:BO:still}
\glll mi·meʔ ṭʼo ʔoʔ ʔacaʔ mukʰúyʔtahqaw\\
    mi·meṭ =ʔṭʼo ʔoc ʔaca·c mukʰuyʔtahqa-w\\
    at.that.time \gl{emph} still person burn.\gl{pl}-\gl{dvs}\\
\glt `At that time they still cremated people.' ({KT 45:11.1})\\
\z
}%

{%
\ea 
\label{ex:BO:notyet1}
\glll ʔoʔ da·bícʰqaw ʔtʰin\\
    ʔoc da·bicʰqa-w ʔ=tʰin\\
    still a.few.go.up-\gl{dvs} \gl{asrt}=\gl{neg}\\
\glt `They had not yet started to go.' ({KT 5:6.2})\\
\z
}%

{%
\ea 
\label{ex:BO:notyet2}
\glll ʔóʔ miʔdaqʰanʔ ma·du·tʰehni\\
    ʔoc mi-ʔdaqʰanʔ ma·duc-tʰ-a-e·=hni\\
    still 2-husband.\gl{nom} arrive-\gl{neg}-\gl{vis.ipfv}-\gl{nfnl}=Q\\
\glt `Didn’t your husband return yet?' ({KT 14:7.3})\\
\z
}%

Expressing `not yet' as `still not' is common cross-linguistically, see e.g.\ \citet{VeselinovaDevos2021} for Bantu languages.

\newpage
\subsection{Negation in non-declaratives}\label{sec:BO:proh}

Positive imperatives are formed by suffixing the verb with \textit{-i} (for a singular addressee) or \textit{-me} (plural addressee, and polite imperative used with in-laws). A glottal stop is often added at the end of the command, as in \ref{ex:kash:12}.%
\footnote{%
See \citet[335--336]{Buckley1994} and \citet{Buckley2015} for the distribution of the final glottal stop, here given the separate gloss \gl{cmnd}, as in `command'.
} %


{%
\ea\label{ex:kash:12}
\glll talá·mecʼi mi· tow ʔ\\
    htala·mecʼ-i mi· =tow =ʔ\\
    climb.down-\gl{imp} there =from =\gl{cmnd}\\
\glt `Climb down from there!' ({KT 21:5.4})\\
\z
}%

Negative imperatives are expressed by means of the suffixes \textit{-tʰu} (for a singular addressee) or \textit{-tʰume}, plus the glottal stop at the end of the clause.
Examples are in (\ref{ex:BO:proh1}) and (\ref{ex:BO:proh2}). In the descriptions of Kashaya morphophonology in \citet{Oswalt1961} and \citet{Buckley1994}, the negative suffix \textit{-tʰ} is treated as having a special allomorph \textit{-tʰu}, which only occurs before the Singular Imperative \textit{-i} (which does not surface) and the Plural Imperative \textit{-me}. I opt for a less abstract description here and treat them as dedicated prohibitive suffixes.

{%
\ea 
\label{ex:BO:proh1}
\glll qaʔcʼáʔtʰuʔ\\
    qaʔcʼaṭ-tʰu=ʔ\\
    cry-\gl{proh}=\gl{cmnd}\\
\glt `Don't cry!' ({KT 9:7.10})\\
\z
}%

\ea 
\label{ex:BO:proh2}
\glll bahcil hayómtʰumeʔ \\
    bahcil hayom-tʰume=ʔ\\
    far many.go.here.and.there-\gl{proh.pl}=\gl{cmnd}\\
\glt [Their mother told them:] `Don't go far!' ({KT 7:4.4})\\
\z

The negative clitic \textit{=tʰin} does not occur in prohibitives.

\subsection{Negation in stative predications}\label{sec:BO:nonverb} 
Stative verbs, such as `know' in \REF{ex:BO:1}, are negated using the same two options as other verbs in the language. In this section, I describe the negation of nominal and existential predications.

\subsubsection{Nominal and adjectival predication}\label{sec:BO:nomadj}
Nominal and adjectival predication is expressed without any lexical verb. Usually, the copula- or auxiliary-like Assertive element \textit{=ʔ} is present in the clause. The Assertive can act as a host to inflectional morphology that would otherwise have been borne by a full lexical verb (such as evidentials, but typically only the `Non-Final Verb' suffix \textit{e·}). Negation is expressed by attaching the negative clitic \textit{=tʰin} directly to the relevant noun phrase.  Nominal predication illustrating proper inclusion (`X is a Y') is seen in examples (\ref{ex:BO:nompred1}) and (\ref{ex:BO:nompred2}). The negative suffix \textit{-tʰ} is not used in non-verbal clauses.


{%
\ea 
\label{ex:BO:nompred1}
\glll {ʔama· dunawá·du hqal} e· maʔu dúwi\\
    {ʔama· dunawa·du hqal} =ʔ-a-e· maʔu duwi\\
    cheater =\gl{asrt}-\gl{vis.ipfv}-\gl{nfnl} this.\gl{nom} coyote\\
\glt `This Coyote is a cheater.' ({KT 4:8.2})\\
\z
}%

{%
\ea 
\label{ex:BO:nompred2}
\glll ʔa· ṭʼo walé·pu tʰin e muʔnati ʔa ku·yi cadu ʔé· wale·pu\\
    ʔa· =ʔṭʼo wale·pu =tʰin ʔ-a-e· muʔnati ʔa· ku·yi cad-u ʔ-a-e· wale·pu\\
    \gl{1sg.nom} =\gl{emph} feather.sorcerer =\gl{neg} \gl{asrt}-\gl{vis.ipfv}-\gl{nfnl} but \gl{1sg.nom} once see-\gl{dvs} \gl{asrt}-\gl{vis.ipfv}-\gl{nfnl} feather.sorcerer\\
\glt `I am not a walépu but I once saw one.' ({KT 40:1.2})\\
\z
}

The same structure is used in contexts of negated equation (`X is not Y'), as in (\ref{ex:BO:negeq}), and negated adjectival predication, as in (\ref{ex:BO:adj}).

\ea 
\label{ex:BO:negeq}
\glll meʔe tʰín e· mu\\
      meʔe =tʰin ʔ-a-e· mu\\
      2.father.\gl{nom} =\gl{neg} \gl{asrt}-\gl{vis.ipfv}-\gl{nfnl} that.\gl{nom}\\
\glt `He is not your father.' ({KT 13:2.3})\\
\z

\ea 
\label{ex:BO:adj}
\begin{xlist}
\ex
\glll  pišun é· mu\\
      pišun =\O-a-é· mu\\
      strange =\gl{asrt-vis.ipfv-nfnl} that\\
\glt `It is strange.'
\ex
\glll  pišúnʔ tʰin e· mu\\
      pišúnʔ =tʰin =\O-a-e· mu\\
      strange =\gl{neg} =\gl{asrt-vis.ipfv-nfnl} that\\
\glt `It is not strange.' ({Kashaya dictionary, entry \textit{pišunʔ}})\\
\end{xlist}
\z


\subsubsection{Locative and existential negation: \textit{cʰow} `not exist'}
\label{sec:BO:chow}

Location and existence can be predicated by means of the semantically general verb \textit{ʔi} `be, remain', as illustrated in (\ref{ex:BO:locpos}), or by more specific positional verbs (`sit', `stand', etc.).%
\footnote{%
The textual corpus contains no clear instances of purely locational (`the cat is on the mat') or existential (`there are ghosts') predications in the affirmative; but their negated counterparts, as described in this section, are amply attested.
} %

{%
\ea 
\label{ex:BO:locpos}
\glll ʔahca· hiʔbayá ʔiw\\
 ʔahcaw hiʔbaya ʔi-w\\
 in.house man be-\gl{dvs}\\
\glt `There is a man in the house.' ({Kashaya Dictionary, entry \textit{ʔi}})\\
\z
}%

\noindent A special verb is used for the negation of locative and existential clauses: \textit{cʰow} (from the root /hcʰo·/), which I gloss as `not exist'.
Example (\ref{ex:BO:loc}) shows a negated locative predication, and (\ref{ex:BO:exist2}) a negated existential. Dedicated negative existentials such as Kashaya \textit{cʰow} are common cross-linguistically (see \citealt{Veselinova2013neg.exist}).

\ea 
\label{ex:BO:loc}
\glll mi·meʔ miyá·konʔ ʔaca· cʰow\\
    mi·meṭ miya·-ko·nʔ ʔaca· hcʰo·-w\\
    at.that.time 3-sister's.husband.\gl{nom} at.home not.exist-\gl{dvs}\\
\glt `At that time, his sister’s husband was not at home.' ({KT 12:11.1})\\ % also in KT 21:11.2
\z

\ea 
\label{ex:BO:exist2}
\glll mi·meʔ ṭʼo ʔacaʔ cʰow\\
 mi·meṭ =ʔṭʼo ʔaca·c hcʰo·-w\\
 at.that.time =\gl{emph} person not.exist-\gl{dvs}\\
\glt `At that time, there were no human beings.' ({KT 1:3.1})\\
\z


Possessive predication is expressed by means of the proprietive postposition \textit{=qʼo} `with, having'. There are two competing options for the negative equivalent:  the negative existential verb \textit{cʰow} `not exist', or the privative postposition \textit{=cʰoʔ} `without' (phonemic /cʰot/). A minimal triplet is in (\ref{ex:BO:poss}). See also \sectref{sec:BO:deriv}.

\ea 
\label{ex:BO:poss}
\begin{xlist}
  \ex
\glll pé·su qʼo ʔe· ya\\
    pe·su =qʼo =ʔ-a-e· ya\\
    money =with =\gl{asrt-vis.ipfv-nfnl} \gl{1pl.nom}\\
\glt    `We have money.'
\ex
\glll pe·su cʰowé· yaʔkʰe\\
  pe·su hcʰo·-wa-e· yaʔkʰe\\
  money not.exist-\gl{vis.ipfv}-\gl{nfnl} \gl{1pl.obl}\\
\glt `We don't have money.'
  \ex % Dict entry ʔcʰot
\glll pé·su cʰotʼe· {} ya\\
    pe·su =cʰot =ʔ-a-e· ya\\
    money without =\gl{asrt}-\gl{vis.ipfv}-\gl{nfnl} \gl{1pl}.\gl{nom}\\
\glt `We don't have money.' ({Kashaya Dictionary, entry \textit{qʼo}})\\
\end{xlist}
\z

The privative \textit{=cʰoʔ} `without' is probably related to the root of `not exist', /hcʰo·/. Some other forms that seem to be derivationally related to /hcʰo·/ are in \tabref{tab:BO:not.exist}. The verb root /hcʰoc/ in (b--d) is probably derived from the negative existential /hcʰo·/ `not exist' by suffixation of the Semelfactive \textit{-c}, literally `start to not exist'.

\begin{table} 
\caption{Expressions related to /hcʰo·/ `not exist'}
\label{tab:BO:not.exist}
 \begin{tabularx}{\textwidth}{llll}
  \lsptoprule
&Expression & Analysis & Meaning \\
  \midrule
a. &  \itshape cʰo·nati & /hcʰo·-nati/ [not.exist-\gl{conc.ss}]& `barely'\\
b. & \itshape bahcʰobáhcʰoʔ & /ba-hcʰoc/ [\gl{mouth}-disappear] & `wear away with mouth'\\
c. & \itshape cʰoyíʔ & /hcʰoc-icʼ/ [disappear-\gl{refl}] & `die'\\
d. & \itshape cʰoyicʼqa & /hcʰoyicʼ-hqa/ [die-\gl{caus}] & `kill' \\
  \lspbottomrule
\end{tabularx}
\end{table}

\subsection{Negation in non-main clauses}
\label{sec:BO:nonmain}

\subsubsection{Complement clauses}
\label{sec:BO:SubClauses}
Complement clauses of verbs such as `want', `know' and `say' differ from independent main clauses in their strictly verb-final constituent order (this is a feature they share with other types of subordinate clauses in Kashaya) and in the limited morphological possibilities of the subordinate verb (\citealt{Olsson2010}). The subordinate verb in such clauses is always suffixed with the Default Verbal Suffix, or, if the clause has future or modal semantics, a suffix such as the Future \textit{-ʔkʰe} or (more rarely) Conditional \textit{-iʔba}. As pointed out in \sectref{sec:BO:distr}, the Default Verbal Suffix and the Future suffix are incompatible with the negative suffix \textit{-tʰ} and only allow negation by means of the clitic \textit{=tʰin}. The Default Verbal Suffix and the Future are the overwhelmingly most common suffixes on verbs in complement clauses, so it follows that the most common negation strategy in complement clauses is the clitic \textit{=tʰin}, as in (\ref{ex:BO:compl1}).

\ea 
\label{ex:BO:compl1}
\glll ʔul {} {duwe lébaṭʰemʔ} qowícʼkʰe tʰin {} šiyiʔ\\
    ʔul {} {duwe lebaṭʰemʔ} hqowic-ʔkʰe =tʰin {} hšiyicʼ-ʔ\\
    then {[} midnight return-\gl{fut} =\gl{neg} {]} say.about.self-\gl{dvs}\\
\glt `He said that he would not go back home until midnight.' ({KT 3:12.4})\\
\z

The negative suffix \textit{-tʰ} is used if the subordinate verb is marked by the Conditional \textit{-iʔba}, since this suffix permits suffixal negation (\ref{ex:BO:cond}). This shows that the preference for \textit{=tʰin} in complement clauses is not a direct consequence of their status as non-main clauses, but depends on which suffix marks the subordinate verb.

\ea 
\label{ex:BO:cond}
\glll [\ldots] dú·ciʔ {} mul ʔacaʔ yaʔ heʔen dacʰacíʔtʰiʔba {} [\ldots]\\
    {} ʔdu·cicʼ-ʔ {} mul ʔaca·c =yac heʔen dacʰac-icʼ-tʰ-iʔba {} {}\\
    {} know-\gl{dvs} {[} that person =\gl{agt} how steal-\gl{refl-neg-cond} {]} {}\\
    \glt `[Then the boss had the Indians guard it,] knowing that Indians wouldn’t steal it [\ldots].' ({KT 64:15.3})
\z

\subsubsection{Switch-reference clauses}
\label{sec:BO:SR}
~\\[-2em]

\noindent
\parbox{\textwidth}{%some weird bug means that the paragraph following the preceding section title is not justified
Kashaya has an extensive system of switch-reference suffixes that are used to form clause chains (`X happened, then Y happened, and then Z happened') or  adverbial-like clauses (`if X happens, then Y happens'). Descriptions of the Kashaya switch-reference system are in \citet{Oswalt1961,Oswalt1983} and \citet[35--44]{Olsson2010}. Switch-reference clauses can be negated by the suffix \textit{-tʰ}, as in (\ref{ex:BO:SR1}) and (\ref{ex:BO:SR2}), or by the clitic \textit{=tʰin}, as in (\ref{ex:BO:SR3}). The counts in \tabref{tab:BO:freqs} show that the suffixal option is somewhat more common in the corpus.}

% \mm{@LSP: for some reason this paragraph is not right-justified, can you fix it?}

{%
\ea 
\label{ex:BO:SR1}
\glll ma moʔóʔtatʰipʰila ʔa· ʔé· ma·dal ci·cʼínʔkʰe\\
     ma moʔoʔta-tʰ-pʰila ʔa· =ʔ-a-e· ma·dal ci·cʼid-ʔkʰe\\
     \gl{2sg.nom} whip.with.stick-\gl{neg-irr.ds} \gl{1sg.nom} =\gl{asrt-vis.ipfv-nfnl} \gl{3sg.f.acc} do-\gl{fut}\\
\glt `If you don’t whip her, I’ll do it.' ({KT 29:4.1})\\
\z
}%

{%
\ea 
\label{ex:BO:SR2}
\glll {\ldots} maʔa cóʔdoʔ tiya·col hqayíʔtʰeti\\
    {} maʔa coʔdoq-ʔ tiya·col hqayicʼ-tʰ-eti\\
    {} food give.several-\gl{dvs} \gl{3pl.refl.acc} ask.for.\gl{pl}-\gl{neg-conc.ds}\\
\glt `[\ldots] and [they] gave them food even though they didn’t ask for it.' ({KT 57:10.4})\\
\z
}%

\ea 
\label{ex:BO:SR3}
\glll mulido mul ma·dúʔli tʰin kʼeša·cʼiwaʔ [\ldots]\\
    mulido mul ma·duc-ʔli =tʰin kʼe·š-acʼ-iwacʼ-ʔ {}\\
    it.is.said that one.arrive-\gl{seq.ds} =\gl{neg} search-along.\gl{pl}-\gl{distr.pl-dvs} {}\\
\glt `When he did not return home, they searched far and wide [\ldots].' ({KT 52:2.2})\\
\z

I have not been able to find any semantic difference that explains the choice between suffixal and clitic negation of switch-reference clauses. The only exception is the use of the negative clitic \textit{=tʰin} in (\ref{ex:BO:SRclit}), which presumably is motivated by the clause being in the scope of replacive focus. The intended meaning is `not because I happened to feel like it', whereas the use of the negative suffix \textit{-tʰ} probably would correspond to the reading `because I didn't feel like it'.

\ea 
\label{ex:BO:SRclit}
\glll mul cáhnoʔte· ʔa bíhqatol nihín tʼalahqaba tʰin \\
    mul cahno·d-te-e· ʔa bihqatol nihin tʼala-hqa-ba =tʰin \\
    that talk-\gl{prfm.intv-nfnl} \gl{1sg.nom} at.this.time to.self start.to.feel-\gl{caus-seq.ss} =\gl{neg} \\
\glt `I must say this at this time not because I just happen to feel like it.' (But because I've been told by God.) ({KT 78:6.4})\\
\z

\noindent The use of the negative clitic on constituents within narrow scope is further discussed in \sectref{sec:BO:scope}.

\subsubsection{Negative purpose: `lest'-clauses}
\label{sec:BO:NegPurp}
Purpose clauses are signalled by the Purposive suffix \textit{-ti} on the verb (\citealt[39]{Olsson2010}). Only one example of a clause expressing negative purpose has been found in the corpus, and it is formed by means of the standard Purposive \textit{-ti} followed by the negative clitic \textit{=tʰin} (\ref{ex:BO:lest}).

\ea 
\label{ex:BO:lest}
\glll [\ldots] hiʔbayá ʔel ma·caʔ tiyá·coʔkʰe hisʼu· ʔel ʔul ṭʼi· doʔqʼóʔdiwaʔ ma·cʼicʰqati tʰin \\
    {} hiʔbaya =ʔel ma·cac tiya·coʔkʰe hisʼu· =ʔel ʔul ṭʼi· doʔqʼoʔdi-wacʼ-ʔ ma·c'ic-qa-ti =tʰin \\
    {} man =\gl{acc} \gl{3pl.nom} \gl{3pl.refl.poss} arrow =\gl{acc} already all prepare-\gl{dur.pl}-\gl{dvs} one.arrive.\gl{pl.agt}-\gl{caus-purp} =\gl{neg} \\
\glt `[Then] the menfolk prepared their arrows so as to prevent [the giant] from coming in.' ({KT 21:10.4})\\
\z


\subsection{Negative lexicalisations}

No clear examples of inherently negative lexical items have been found. Possible exceptions (admittedly more grammatical than lexical) are the verb \textit{cʰow} `not exist' and the postposition \textit{=cʰoʔ} `without', as discussed in \sectref{sec:BO:chow}. Another candidate is \textit{naʔbaʔ} `not know how' (root /naʔbac/), but the few attestations of this verb make it impossible to decide whether it is truly inherently negative. (See also \sectref{sec:BO:replies} for the answers `no' and `I don't know'.)

\subsection{Other clausal negation constructions: The inability construction}
\label{sec:BO:inabil}

Ability -- `can, be able to' -- is mainly expressed by suffixing the verb with the Conditional \mbox{\textit{-iʔba}} or Future \textit{-ʔkʰe} \citep[251, 257]{Oswalt1961}. Interestingly, inability is expressed not by simply negating such expressions, but by negating them and adding an interrogative word, typically \textit{heʔen} `how', to the clause. The presence of the interrogative word makes the structure a non-compositional counterpart of positive expressions of ability (which occur without an interrogative word), and the structure must therefore be considered a separate construction.
The inability construction is very frequent in the corpus; examples are in (\ref{ex:BO:inabil1}) and (\ref{ex:BO:inabil2}). In fact, instances of this construction are much more frequent than expressions of (positive) ability, which often seems to be expressed implicitly (i.e.\ `I see' rather than `I can see').

\ea 
\label{ex:BO:inabil1}
\glll heʔen mu·kinʔ hiʔda daʔtʼáʔkʰe tʰin\\
    heʔen mu·kinʔ hiʔda daʔtʼa-ʔkʰe =tʰin\\
    how \gl{3sg.m.nom} path find-\gl{fut} \gl{neg}\\
\glt `He couldn’t find the trail.' ({KT 27:4.6})\\
\z

\ea 
\label{ex:BO:inabil2}
\glll [\ldots] mensʼin ya heʔen bayatáʔkʰe tʰin ma·caʔkʰe ya kúʔmul sʼiʔí ʔnati\\
    {} mensʼin  ya heʔen bayataq-ʔkʰe =tʰin ma·caʔkʰe ya kuʔmul sʼiʔi ʔ-nati\\
    {} and \gl{1pl.nom} how understand.\gl{pl-fut} =\gl{neg} \gl{3pl.obl} \gl{1pl.nom} same flesh \gl{asrt}-\gl{conc.ss}\\
\glt `[Other people have different languages] and we can’t understand theirs even though we are of one flesh.' ({KT 2:8.7})\\
\z

Interrogative expressions such as \textit{ciba·} `who' function as negative indefinite pronouns in negated clauses (`nobody' etc.; see \sectref{sec:BO:negindef} below), so \textit{heʔen} `how' should perhaps be interpreted as meaning something like `in no way' in these contexts. But \textit{heʔen} `how' is usually omitted if some other interrogative expression is present, as in (\ref{ex:BO:inabil3}) and (\ref{ex:BO:inabil4}), so the generalisation seems to be that the expression of inability requires some interrogative word to be present, with the default being \textit{heʔen} `how'.

\ea 
\label{ex:BO:inabil3}
\glll hiʔdi ʔahqʰa daʔtʼáʔkʰe tʰin\\
    hiʔdi ʔahqʰa daʔtʼa-ʔkʰe =tʰin\\
    where water find-\gl{fut} =\gl{neg}\\
\glt `He couldn’t find water anywhere.' ({KT 2:9.4})\\
\z

\ea 
\label{ex:BO:inabil4}
\glll muʔnati ciba· mul cácʼkʰe tʰin\\
    muʔnati ciba· mul cacʼ-ʔkʰe =tʰin\\
    but who that see.\gl{pl.agt-fut} \gl{neg}\\
\glt `But nobody can see that.' ({KT 37:7.3})\\
\z


\section{Non-clausal negation}

\subsection{Negative replies}
\label{sec:BO:replies}


One-word replies to polar questions are: \textit{hu·ʔ} `yes', \textit{da·} `no' (and \textit{huhu·} `I don't know'). If the question is negated, as in (\ref{ex:BO:da.a}), the reply \textit{da·} `no' disagrees with the content of the question, not with its polarity: in (\ref{ex:BO:da.b}) Cricket Woman is replying that her husband did not return home.


{%
\ea 
\label{ex:BO:da}
partly repeated from (\ref{ex:BO:notyet2})
    \begin{xlist}
    \ex\label{ex:BO:da.a}
        \glll \llap{``}ʔóʔ miʔdaqʰanʔ ma·du·tʰehni'' hcedu\\
            ʔoc mi-ʔdaqʰanʔ ma·duc-tʰ-a-e·=hni hced-u\\
            still \gl{2}-husband.\gl{nom} arrive-\gl{neg}-\gl{vis.ipfv}-\gl{nfnl=q} say-\gl{dvs}\\
        \glt `[He] asked, “Didn’t your husband return yet?''\,'
    \ex\label{ex:BO:da.b}
        \glll \llap{``}dá·'' cedu ṭʰoʔo·koyʔmenʔ\\
             da· hced-u ṭʰoʔo·koyʔ-menʔ\\
             no say-\gl{dvs} cricket-\gl{female.name}\\
        \glt `“No,” said Cricket Woman.' ({KT 14:7.3--4})\\
    \end{xlist}
\z
}%

Positive replies to negative questions use \textit{hu·ʔ} `yes'. In (\ref{ex:BO:huu}), the verb from the question is repeated, but this is not an obligatory feature of positive replies to negative questions (see e.g. KT6:4.11). 

\ea 
\label{ex:BO:huu}
\begin{xlist}
\ex
\glll {[\ldots] } \llap{``}ma ṭʼo šo·me qʼoʔo dú·ciʔ tʰin ehni men ʔbakʰe'' [\ldots]\\ %HC search ldots in this chapter and see if I can make them look better
    {} ma =ʔṭʼo šo·me qʼoʔo ʔdu·cicʼ-ʔ =tʰin ʔ-a-e·=hni men =bakʰe {}\\
    {} \gl{2sg.nom} =\gl{emph} younger.sister.\gl{voc} song know-\gl{dvs} =\gl{neg} \gl{asrt-vis.ipfv-nfnl=q} thus for\\
\glt `[Then his mother said to my grandmother,] ``Sister, you don’t know a song for that, do you?”
\ex
\glll mulido \llap{``}hú·ʔ dú·cicʼem ṭa ʔa'' nihcedu\\
    mulido hu·ʔ ʔdu·cicʼ-em =ṭa ʔa nihced-u\\
    it.is.said yes know-\gl{resp} =\gl{foc} \gl{1sg.nom} say-\gl{dvs}\\
\glt `“Yes, I know one,” she replied.' ({KT 41:11.2})\\
\end{xlist}
\z


\subsection{Negative indefinites}
\label{sec:BO:negindef}

Like in many other languages, interrogative expressions such as \textit{ciba·} `who' are used as negative indefinites in negated clauses, as in (\ref{ex:BO:negindef1}--\ref{ex:BO:negindef3}). There is no ambiguity with content questions, because in content questions the Interrogative \textit{-wa} is present in the clause \citep[286]{Oswalt1961}.

{%
\ea 
\label{ex:BO:negindef1}
\glll ciba· duyáʔtaw ʔtʰin\\
    ciba· duyaʔta-w ʔ=tʰin\\
    who help-\gl{dvs} \gl{asrt=neg}\\
\glt `Nobody helped him.' ({KT 2:4.2})\\
\z
}%


\ea 
\label{ex:BO:hide}
\glll cibal ya qʼawá·cʼi·tʰela\\
    ciba·l ya qʼawa·cʼic-tʰ-ela\\
    who.\gl{acc} \gl{1pl.nom} hide.\gl{pl}-\gl{neg}-\gl{prfm.ipfv}\\
\glt `We haven’t hid anybody.' ({KT 16:12.6})\\
\z

{%
\ea 
\label{ex:BO:negindef3}
\glll mulido baqʼo maci·biʔ tʰin\\
    mulido baqʼo ma·c-ibic-ʔ =tʰin\\
    it.is.said what burn-\gl{incep}-\gl{dvs} =\gl{neg}\\
\glt `Nothing started to burn.' ({KT 16:14.3})\\
\z
}

As seen in (\ref{ex:BO:negindef4}), interrogatives have the negative indefinite meaning in clauses with \textit{cʰow} `not exist' too, which shows that this verb must be regarded as inherently negative.

\ea 
\label{ex:BO:negindef4}
\glll hḿ··· ʔacaʔ mihšewé· [\ldots]. cʰo·wé· cibá· [\ldots]\\
    hm ʔaca·c mihše-wa-e· {} hcʰo·-wa-e· ciba· {}\\
    hm person smell-\gl{vis.ipfv}-\gl{nfnl} {} not.exist-\gl{vis.ipfv}-\gl{nfnl} who {}\\
\glt `\,``Hmm, I smell a person'' [said Eagle]. ``There is no one'' [said Old Lady Kestrel].' ({H 4:6.5})\\
\z

Unfortunately, there is no data showing whether the negative indefinite use of interrogatives extends across clause boundaries, i.e.\ whether one can say (the equivalent of) `I don't think \textsc{who} helped him' to mean `I don't think anybody helped him' (so called indirect negation).

A remarkable feature of Kashaya is that the formation of positive indefinite pronouns also involves negation: these are formed by adding the negative clitic directly to the question words, as described in \sectref{sec:BO:indef}.

\subsection{Negative derivation and case marking}
\label{sec:BO:deriv}

The only known item that can be said to express negative meanings on the phrasal level is \textit{cʰoʔ} `without' (\sectref{sec:BO:chow}), as in (\ref{ex:BO:dirty}). This utterance is an example of standard non-verbal predication (as discussed in \sectref{sec:BO:nonverb} above), with the postpositional phrase \textit{sʼaʔsʼa cʰoʔ} `clean' (lit. `without dirt') functioning as the predicate. All examples of such expressions that I found are used predicatively (either as non-verbal predicates or as depictive secondary predicates), so it is not known if these expressions can be used attributively, as adnominal modifiers (`a clean cloth').

\ea 
\label{ex:BO:dirty}
\glll mu ṭʼo sʼaʔsʼa cʰoʔ\\
    mu =ʔṭʼo sʼaʔsʼa =cʰot\\
    that =\gl{emph} dirt without\\
\glt `It was not dirty.' ({KT 5:4.5})\\
\z

\newpage
Some expressions containing \textit{cʰoʔ} `without' are in \tabref{tab:BO:deriv}, taken from \citet{Oswalt}. It is not known to what degree these are established idioms or nonce creations.

\begin{table}
\caption{Expressions with \textit{cʰoʔ} `without'}
\label{tab:BO:deriv}
 \begin{tabular}{llll}
  \lsptoprule
  Expression & gloss & meaning \\
  \midrule
  \itshape maʔa cʰoʔ & food without & `hungry' & \\
  \itshape šimamo cʰoʔ & inner.ear without & `flirtatious man' \\
  \itshape ciliʔ cʰoʔ & be.worth without & `cheap' \\
  \itshape cahno cʰóʔ & speech without & `quiet' \\
  \itshape ʔama cʰíyaʔ cʰoʔ & thing be.afraid without & `brave, fearless' \\
  \lspbottomrule
 \end{tabular}
\end{table}

\section{Other aspects of negation}

\subsection{The scope of negation}
\label{sec:BO:scope}

There appear to be several options for narrowing down the scope of negation to a specific constituent. The most frequent option is to mark the constituent in the scope of negation directly, by means of the negative clitic \textit{=tʰin}. In (\ref{ex:BO:negfoc1}), the expression \textit{ʔacaʔ} `person, people' is negated directly, as it is in the scope of contrastive negation.

{%
\ea 
\label{ex:BO:negfoc1}
\glll mensʼiba miqʰama·tow ʔṭʼo mul ʔacaʔ tʰin nihciʔ šiʔbaši nihciʔ\\
    mensʼiba miqʰama·tow =ʔṭʼo mul ʔaca·c =tʰin nihcicʼ-ʔ šiʔbaši nihcicʼ-ʔ\\
    and.then after.that =\gl{emph} that person =\gl{neg} say.\gl{pl.agt-dvs} animal say.\gl{pl.agt-dvs}\\
\glt `Thereafter they weren’t called people; they were called animals.' ({KT 10:22.8})\\
\z
}%

Second, focus particles (usually the Emphatic focus particle \textit{=ʔṭʼo}) can be used, as in \REF{ex:BO:45}.
    

\ea \label{ex:BO:45}
\glll da·qáʔ tʰin ku ṭʼo\\
    da·qacʼ-ʔ =tʰin ku =ʔṭʼo\\
    want-\gl{dvs} =\gl{neg} one =\gl{emph}\\
\glt `He did not want just one.' (i.e., he wanted several) ({KT 4:9.4})\\
\z
    
Finally, constituent focus can be expressed with standard clausal negation, as in (\ref{ex:BO:foc3}). Prosodic prominence is perhaps relevant for the interpretation of such utterances.%
\footnote{Oswalt's original recordings are archived in the Survey of California and Other Indian Languages (\url{http://cla.berkeley.edu/collection/11055}), but analysis of the interaction between prosody and negation was outside the scope of this study.}

\ea 
\label{ex:BO:foc3}
\glll [\ldots] ma·caʔ ṭʼo mul ʔacaʔ sʼiʔi bimuyíʔ tʰin, bihše máʔyul bimuyiʔ [\ldots]\\
    {} ma·cac =ʔṭʼo mul ʔaca·c sʼiʔi bimuyicʼ-ʔ =tʰin bihše maʔyul bimuyicʼ-ʔ {}\\
    {} \gl{3pl.nom} \gl{emph} that person flesh eat.\gl{pl}-\gl{dvs} =\gl{neg} deer only eat.\gl{pl}-\gl{dvs}\\
\glt `[\ldots] they didn’t dine on human flesh; they only ate meat [\ldots]' ({KT 6:2.3})\\
\z

See also the discussion of negative scope in complex clauses in \sectref{sec:BO:SRscope}.

\subsection{Negative polarity}

No idiomatic negative polarity items of the type \textit{lift a finger} are attested in the textual corpus or in the Kashaya Dictionary.
The use of interrogatives as negative indefinites under negation was noted in \sectref{sec:BO:negindef}.

\subsection{Marking of NPs in the scope of negation}

\label{sec:BO:flag}
Like the other languages of the Pomoan family, Kashaya has a case system
based on semantic alignment (see e.g.\ \citealt{Mithun1991}). The Nominative case clitic \textit{=ʔem} marks NPs expressing agent-like participants (e.g.\ the sole argument of an agentive verb such as `dance') and the Accusative \textit{=ʔel} marks patient-like participants (e.g.\ the sole argument of a patientive verb such as `die'). Certain expressions, such as personal pronouns, kinship terms, and clausal nominalisations, are obligatorily marked for case (marked by suffixes rather than clitics). Lexical NPs seem to be case-marked under discourse conditions that are not yet well understood, but it seems to be mainly referential, definite NPs that are marked for Nominative and Accusative case. These principles seem to be at work in negated clauses too, so a definite NP such as `the bear' in (\ref{ex:BO:bear}) is overtly marked by case, whereas a non-referential NP such as `water' in (\ref{ex:BO:inabil3}) above is unmarked for case.

\ea 
% ending of a story in which a bear falls into a river.
\label{ex:BO:bear}
\glll miqʰamá·to· buṭaqá ʔel caʔ tʰin\\
    miqʰama·tow buṭaqa =ʔel cacʼ-ʔ =tʰin\\
    after.that bear =\gl{acc} see.\gl{pl.agt-dvs} =\gl{neg}\\
\glt `After that the bear was never seen again.' ({KT 5:19.5})\\
\z

At present, then, there is no evidence suggesting that case marking is affected by negation, but this conclusion must be regarded as preliminary pending more work on case marking in Kashaya.

\subsection{Reinforcing negation}
Reinforcing or minimising expressions such as `(not even) one' or `(not even) a little' are formed with the element \textit{-nati} (usually attached to the Assertive \textit{ʔ}), as in (\ref{ex:BO:even1}) and (\ref{ex:BO:even2}). (\textit{-nati} is also a concessive marker meaning `although'.)

{%
\ea 
\label{ex:BO:even1}
\glll kú ʔnati dacew ʔtʰin\\
    ku ʔ-nati dace-w ʔ=tʰin\\
    one \gl{asrt}-even catch-\gl{dvs} \gl{asrt=neg}\\
\glt `They didn’t catch even one.' ({KT 59:12.12})\\
\z
}%

{%
\ea 
\label{ex:BO:even2}
\glll qawí ʔnati maʔal miʔkʰe míhṭʰe di·cʼíʔtʰuʔ\\
    qawi ʔ-nati maʔal miʔkʰe mi-hṭʰe di·cʼ-id-tʰu =ʔ\\
    little \gl{asrt}-even this.\gl{acc} \gl{2sg.poss} 2-mother tell-\gl{dur}-\gl{neg} =\gl{imp}\\
\glt `Don’t breathe a word to your mother.' ({KT 29:14.2})\\
\z
}%

\subsection{Negation, coordination and complex clauses}
\label{sec:BO:SRscope}
The main strategy for coordinating complex clauses is the system of switch-reference, as discussed in \sectref{sec:BO:SR}. I already discussed negation of the non-main clauses in such constructions; here I will add some information about the scope that negation takes when it is expressed in the main clause. There are two possibilities for an utterance containing a switch-reference clause and a main clause in which the main verb is negated. Either the main clause negation has scope only over the main clause, as in (\ref{ex:BO:toy.a}), or over both clauses, as in (\ref{ex:BO:toy.b}).

{%
\ea Having done A, I \textsc{neg} did B.
\label{ex:BO:toy}
\begin{xlist}
\ex `I did A, and I didn't do B.'\label{ex:BO:toy.a}
\ex `I didn't do A, and I didn't do B.'\label{ex:BO:toy.b}
\end{xlist}
\z
}

Both of these options are frequent in the corpus.
Examples of negation with scope over only the main clause, as in (\ref{ex:BO:toy.a}) above, are in (\ref{ex:BO:scope1}) and (\ref{ex:BO:scope2}).

\ea 
\label{ex:BO:scope1}
\glll \llap{``}bihše boʔoba ma·dú·tʰe·'' nihcedu\\
    bihše boʔo-ba ma·duc-tʰ-a-e· nihced-u\\
    deer hunt-\gl{seq.ss} one.arrive-\gl{neg-vis.ipfv-nfnl} say-\gl{dvs}\\
\glt `“Having gone hunting, he didn’t come back,” she said.'\\
(NOT: `He didn't go hunting, and he didn't come back') ({KT 7:3.5})\\
\z

{%
\ea 
\label{ex:BO:scope2}
\glll maʔa ʔel dahólmadun daʔtʼaci·du tʰin men haʔdá· caci·du\\
    maʔa =ʔel dahol-m-ad-in daʔtʼa-cid-u =tʰin men haʔdaw ca-cid-u\\
    food =\gl{acc} search-\gl{ess-dur-sim.ss} find-\gl{dur-dvs} =\gl{neg} thus hungry sit-\gl{dur-dvs}\\
\glt `While searching for food [the child] wouldn’t find it; it sat hungry.'\\
(NOT: `When it didn't search for food, it didn't find it') ({KT 51:3.2})\\
\z
}%

An example of negation with scope over both the main clause and the switch-reference clause, as in (\ref{ex:BO:toy.b}) above, is in (\ref{ex:BO:scope3}).

\ea 
\label{ex:BO:scope3}
\glll mensʼin ʔdú·cicʼe· ʔa be·li ʔa qʰaʔbe cícʼpʰi tʼetʰmaʔkʰe tʰin\\
    mensʼin ʔdu·cicʼ-a-e· ʔa be·li ʔa qʰaʔbe ci·cʼ-pʰi tʼe·tʰma-ʔkʰe =tʰin\\
    and know-\gl{vis.ipfv-nfnl} \gl{1sg.nom} here \gl{1sg.nom} rock become-\gl{irr.ss} one.stand-\gl{fut} =\gl{neg}\\
\glt `I know that I am not going to turn into a rock and stand here.' \\
(NOT: `\ldots that when I have turned into a rock, I won't stand here.') ({KT 78:15.7})\\
\z

Similar cases of flexibility in the interpretation of negative scope have been noted for other languages with switch-reference systems (e.g.\ \citealt[258--259]{Reesink2002}).

\subsection{Miscellaneous aspects of negation}
\subsubsection{Non-negative uses of negation: indefinite pronouns} 
\label{sec:BO:indef}

The standard indefinite pronoun series in Kashaya are formed by adding \textit{=tʰin}, i.e.\ the negative clitic, to the interrogatives, as shown in \tabref{tab:BO:indef}. Thus, \textit{ciba· tʰín} `someone' is literally `who not' and \textit{baqʼo tʰín} `something' is literally `what not'.
Indefinite pronouns that are based on interrogative pronouns are very common cross-linguistically (e.g.\ \citealt{Haspelmath1997}), but the combination of interrogative pronouns with negation to form positive indefinite pronouns has not been documented in the literature,%
\footnote{%
Note that the English expression \textit{and whatnot} means `and similar items, etc.' and is not an indefinite pronoun.
} %
and is indeed a very surprising feature of the language.%


\begin{table}[htb]
\caption{Interrogatives and indefinites}
\label{tab:BO:indef}
 \begin{tabular}{llll}
  \lsptoprule
\multicolumn{2}{l}{Interrogative} & \multicolumn{2}{l}{Indefinite} \\
  \midrule
\itshape cibá· & `who' & \itshape ciba· tʰín & `someone' \\
\itshape baqʼó & `what' & \itshape baqʼo tʰín & `something' \\
\itshape buté· & `when' & \itshape bute· tʰín & `sometimes'\\
\itshape béṭʼbu & `how many, how much' & \itshape béṭʼbu tʰin & `a few, some'\\
\itshape heʔéy & `where' & \itshape heʔéy tʰin & `somewhere' \\
\itshape hiʔdí & `whereabouts' & \itshape hiʔdi tʰín & `somewhere' \\
\itshape heʔén & `how' & \itshape heʔén tʰin & `somehow, in some way' \\
  \lspbottomrule
 \end{tabular}
\end{table}

The standard indefinite series are used in most of the non-negative functions identified in \citet{Haspelmath1997}. For example, \textit{ciba· tʰín} `someone' can be used to express a specific, unknown indefinite referent (\ref{ex:BO:indef0a}), and it can be used in conditionals (\ref{ex:BO:indefcond}).

\ea 
\label{ex:BO:indef0a}
\glll {[\ldots]} to cibá· tʰin da·qaʔ\\
{} to ciba· =tʰin da·qac'-ʔ\\
 {} 1\gl{sg}.\gl{acc} who =\gl{neg} like-\gl{dvs}\\
\glt `[That is when I know that] someone wants me.' ({KT 50:18.5})\\
\z

\ea
\label{ex:BO:indefcond}
\glll cibá· tʰin mito ʔama· mito men hqacínʔpʰila [\ldots]\\
 ciba· =tʰin mito ʔama· mito men hqac-id-pʰila {}\\
 who =\gl{neg} 2\gl{sg}.\gl{acc} thing 2\gl{sg}.\gl{acc} thus ask-\gl{dur}-\gl{ds.irr} {}\\
\glt `If anyone ever asks anything of you, [you should not say, ‘No’]'
\z

\noindent Compare these uses of the positive indefinite pronouns with the uses of bare interrogatives as negative indefinite pronouns in negated clauses, described in \sectref{sec:BO:negindef} above. Using the concept 'who' as illustration, the two uses of interrogatives outside content questions can be summarised as follows:

\eabox{
\begin{tabular}{lll}
\textsc{who}+\textsc{neg} & \textsc{verb} & `Someone \textsc{verb}-ed.'\\
\textsc{who} & \textsc{verb}+\textsc{neg} & `Nobody \textsc{verb}-ed.'\\ 
\end{tabular} 
}

The use of negation to form non-negative indefinites is so surprising that it is worth asking whether the element \textit{tʰin} in these expressions is not simply a case of accidental homophony with the negative clitic \textit{=tʰin}, but I am not aware of any other element of the shape \textit{tʰin}, either in Kashaya or in other closely related Pomoan languages, that could have provided the diachronic source for the indefinite marker.

Note also that \textit{=tʰin} clearly behaves as a clitic even in indefinite expressions, because it detaches from the interrogative word when other elements, such as postpositions, are added to the indefinite expression.
In such complex indefinite expressions, \textit{=tʰin} appears at the right edge of the complex interrogative expression, following the added material. In (\ref{ex:BO:indef1}), the interrogative \textit{heʔé·} `where' is followed by the postposition \textit{tow} `from', which intervenes before \textit{=tʰin}, resulting in a complex indefinite \textit{heʔé· tow ʔtʰin} `from somewhere'. The fact that \textit{=tʰin} behaves like a clitic even in indefinite expressions could be considered to be evidence that it is the `same' item as the negative clitic.%
\footnote{%
Note that the negative clitic attaches to the Assertive \textit{ʔ} (which functions as an auxiliary verb) in complex indefinite expressions. The Assertive occasionally hosts the negative clitic in standard clausal negation too, as in (\ref{ex:BO:notyet1}) and (\ref{ex:BO:even1}) above. I am not aware of the function of the Assertive in these uses. A reviewer suggests that \textit{=ʔtʰin} might simply be an allomorph of \textit{=tʰin} in certain contexts, perhaps historically related to the Assertive.
} %

\ea 
\label{ex:BO:indef1}
\glll mulidom pʰiʔtʼan heʔé· tow ʔtʰin cibá· tʰin buʔsʼuná·yadu\\
 mulidom pʰiʔtʼan heʔey =tow ʔ=tʰin ciba· =tʰin buʔsʼunay-ad-u\\
 it.is.said suddenly where =from \gl{asrt}=\gl{neg} who =\gl{neg} make.kissing.noise-\gl{dur}-\gl{dvs}\\
\glt `Suddenly from somewhere someone was making a kissing noise.' ({KT 24:3.2})\\
\z

Comparison with indefinite pronouns in the closest relatives of Kashaya does not shed any light on the diachronic origin of the positive indefinite pronouns.%
\footnote{%
Southern Pomo uses interrogative words such as \textit{čaʔ:a} `who' followed by the clitic \textit{ʔnat̪i} `but, even' as negative indefinite pronouns, combined with a negated verb (data from Halpern's Southern Pomo Texts; cf.\ also \citealt[259]{Walker2020}); I have found no data on positive indefinite pronouns in this language.
Central Pomo uses bare interrogative pronouns (e.g.\ \textit{bá:} `who'; cf.\ Kashaya \textit{ciba·} `who') as indefinite pronouns in both positive and negative contexts (e.g.\ \citealt[238, ex.\ 58]{Mithun2008passivization}, \citealt[83, ex.\ 15]{Mithun1998})
} %
\citet{Haspelmath1997} mentions one source of positive indefinite pronouns that involves negation, the `dunno' type, i.e.\ expressions that originally meant `I don't know who' (e.g.\ Old Norse \textit{nekkver} `somebody' < \textit{*ne wait ik hwarir} `I don't know who'; \citeyear[131]{Haspelmath1997}). Haspelmath notes that `dunno'-type indefinites seem restricted to European languages, and there is nothing in Kashaya indefinites such as \textit{ciba· tʰín} `someone' suggesting that these have a clausal origin. It is worth noting, however, that utterances combining `don't know' and an interrogative word, as in (\ref{ex:BO:dunno}), are not uncommon in the Kashaya corpus, so it is possible that this type of structure is the diachronic source of the indefinite pronouns. This would have involved placement of the negative clitic after the interrogative word (as in constituent negation, \sectref{sec:BO:scope}) and a meaning change, i.e.\ `don't know where X went' > `X went somewhere'.
Unfortunately there is no data to support this scenario, so the explanation for the use of negation in the indefinite pronouns remains an enigma.

{%
\ea 
\label{ex:BO:dunno}
\glll muʔnati heʔé· wa·du dú·ciʔ tʰin [\ldots]\\
 muʔnati heʔey wa·d-u ʔdu·cicʼ-ʔ =tʰin {}\\
 but where one.go.along-\gl{dvs} know-\gl{dvs} =\gl{neg} {}\\
\glt `But they didn’t know where she went; [she was lost.]' ({KT 51:6.3})\\
\z
}%


\subsubsection{Comparative remarks on negation in Pomoan languages}
The existence of multiple negation strategies in Kashaya (suffix \textit{-tʰ}, clitic \textit{=tʰin}, verb \textit{cʰow}) appears to be shared with the other Pomoan languages. The closest relative, Southern Pomo, shows a very similar pattern with competition between a negative suffix \textit{-t̪ʰ} (used with e.g.\ Conditional verb forms) and a negative clitic \textit{=t̪ʰot̪ʼ} or \textit{=t̪ʰoṭʼ} (used with the Future suffix, and in non-verbal predication; \citealt[246--248]{Walker2020}), plus a negative existential verb \textit{ʔačʰ:ów}, corresponding to Kashaya \textit{cʰow} `not exist' (\citealt[308]{Walker2020}).

Central Pomo, spoken north of Kashaya, shows a slightly more complex situation (\citealt{Mithun1998}). The markers of negation are clearly cognate with the negative morphemes in Kashaya and Southern Pomo, but their functions are different. According to Mithun's description, Central Pomo negation is intertwined with aspect, and distinguishes the perfective negator \textit{čʰów} (cognate with Kashaya \textit{cʰow} `not exist') from the imperfective negator \textit{t̪ʰin} (cognate with Kashaya \textit{=tʰin}). Mithun observes that \textit{čʰów} also functions as a negative existential, and that \textit{t̪ʰin} is used in nominal and adjectival predication, just like the Kashaya cognates.

Negation in Eastern Pomo employs the suffixal element \textit{-kʰúy} (\citealt[105]{McLendon1975}), which is cognate with Kashaya \textit{cʰow}, Southern Pomo \textit{ʔačʰ:ów}, and Central Pomo \textit{čʰów}. These reflect Proto-Pomo \textit{*ʔakʰ·ów} (\citealt[83]{McLendon1973}).
The geographic distribution of the descendants of Proto-Pomo \textit{*ʔakʰ·ów} is restricted to existential negation in the southernmost languages (Kashaya and Southern Pomo), extends to negating verbs in some contexts in Central Pomo, and functions as the standard negator in Eastern Pomo. This aligns neatly with different diachronic stages of the so-called Negative Existential Cycle, which describes the development of standard negators from negative existentials (\citealt{Croft1991}; \citealt{Veselinova2014}).

The other Pomoan languages for which data are available appear to feature double-marked negation, combining negative verb forms with pre-verbal negative particles (\citealt{Moshinsky1974}; \citealt{OConnor1987}). Untangling the history of the diverse negation strategies of the Pomoan family would be a rewarding comparative task.

The fascinating `what-not' indefinite pronoun series described in \sectref{sec:BO:indef} lack equivalents in other Pomoan languages. The formation of positive indefinite pronouns from interrogatives and the negative marker appears to be a quirk of Kashaya, and a typological rarissimum to boot.


\section{Summary}
\largerpage[1.5]
\begin{table}[b]
\caption{Summary of negation strategies}
\label{tab:BO:last}
\begin{tabularx}{\textwidth}{lQl}
\lsptoprule
Negation strategy & Functions \\
\midrule
Suffix \textit{-tʰ} &  Negating verbs marked with evidential suffixes, Conditional \textit{-iʔba}, etc. & (\ref{sec:BO:distr})\\
\tablevspace
Clitic \textit{=tʰin} &  Negating verbs marked with Future \textit{-ʔkʰe}, Default Verbal Suffix, etc. & (\ref{sec:BO:distr})\\
		&  Negation in nominal \& adjectival predication & (\ref{sec:BO:nomadj})\\
		&  Constituent negation & (\ref{sec:BO:scope})  \\
    & Formation of positive indefinite pronouns & (\ref{sec:BO:indef})\\
\tablevspace
Verb \textit{cʰow} `not exist' &  Existential \& locative negation & (\ref{sec:BO:chow}) \\
\tablevspace
Suffixes \textit{-tʰu} (sg), \textit{-tʰume} (pl) & Prohibitive & (\ref{sec:BO:proh})\\
\lspbottomrule
\end{tabularx}
\end{table}

Kashaya has three main strategies for expressing negation: a verb suffix \textit{-tʰ}, a clitic \textit{=tʰin} and a negative existential verb \textit{cʰow}. \tabref{tab:BO:last} outlines the distribution of the three strategies.

I showed that the three strategies are largely in complementary distribution, with determinants involving both morphological context (e.g.\ verbs with evidentials only taking the negative suffix \textit{-tʰ}; \sectref{sec:BO:distr}) and syntactic context (e.g.\ constituent negation and nominal/adjectival predication selecting the clitic \textit{=tʰin}; \sectref{sec:BO:nomadj}, \sectref{sec:BO:scope}). I noted that some contexts show variation between the suffixal and clitic options (as in the negation of switch-reference clauses, \sectref{sec:BO:SR}), but overall, the patterning of the three strategies is straightforward.

\section*{Acknowledgements}
I was fortunate to be able to learn about Kashaya from a native speaker, the late Anita Silva, during the Field Methods class led by Pam Munro at the LSA 2009 Institute at UC Berkeley. My work with Mrs.\ Silva focused on negation, but was mainly based on elicitation, so I have chosen to base this chapter on textual data instead. I am indebted to Pam Munro, who encouraged me to work on negation and provided very helpful comments on my class paper, and to Gene Buckley, who kindly shared his interlinearised Kashaya Texts with me. Thanks are also due to the two anonymous reviewers and the editors of this volume, who provided helpful comments on this chapter.


\section*{Abbreviations}
\begin{tabularx}{.45\textwidth}{lQ}
1 & first person\\
2 & second person\\
3 &  third person\\
\scshape acc & accusative\\ 
\scshape agt &  agent\\
\scshape asrt  & assertive (auxiliary)\\  
\scshape caus & causative\\ 
\scshape cmnd  & command\\  
\scshape conc  & concessive\\
\scshape cond  & conditional\\
\scshape distr & distributive\\
\scshape ds  & different subject\\
\scshape dur & durative\\
\scshape dvs & default verbal suffix\\
\scshape emph  & emphatic focus\\
\scshape ess & essive\\
\scshape EVID & evidentials\\%capitals are here because it is not a morpheme in itself but there are different evid markers glossed with EVID in this paper. don't take away the capitals.
\end{tabularx}
\begin{tabularx}{.45\textwidth}{lQ}
\scshape f & feminine\\
\scshape foc & focus\\
\scshape fut & future\\
\scshape gnr & general \\
\scshape hrsy & hearsay evidential\\
\scshape imp & imperative\\
\scshape incep & inceptive\\
\scshape intv & intentive\\
\scshape ipfv & imperfective\\
\scshape irr & irrealis\\
\scshape m & masculine\\
\scshape neg & negative\\
\scshape nfnl & non-final verb\\
\scshape nom & nominative\\
\scshape NP & nominal phrase\\
\scshape obl & oblique\\
\scshape prfm & performative\\
\end{tabularx}

\begin{tabularx}{.45\textwidth}{lQ}
\scshape pfv & perfective\\
\scshape pl & plural\\
\scshape poss & possessive\\
\scshape proh & prohibitive\\
\scshape purp & purposive\\
\scshape q & yes/no-question\\
\scshape refl & reflexive\\
\scshape resp & responsive\\
\end{tabularx}
\begin{tabularx}{.45\textwidth}{lQ}
\scshape seq & sequential\\
\scshape sg & 	singular\\
\scshape sim & simultaneous\\
\scshape SR & switch-reference\\
\scshape ss & same subject\\
\scshape V & verb\\
\scshape vis & visual evidential \\
\scshape voc & vocative \\
\end{tabularx}



\printbibliography[heading=subbibliography,notkeyword=this]

\end{document}
